%% APPENDIX B -- the control assumption, the reserved hold-out and the
%% recurrence bound.  Carries the labels app:radius (04_method.tex,
%% 05_results.tex), sec:heldout (03_sample.tex, 05_results.tex,
%% app_J_falsealarm.tex), sec:nullscatter (05_results.tex) and eq:excl.
%% r30 (write-apx): the assumption that a star behaves like a control
%% position is stated HERE once; Appendix D keeps the estimator and the
%% resampling and points here.  The 7 m-array excess: two sentences, its only
%% place.  Fig. B1 is deleted (the beta Pictoris control is in Sec. 4).
%% Register row 2: the inference "a calibration residual cannot reach the
%% trigger, so it is excluded" is withdrawn; a small continuum-coupled
%% residual can change the RATE without making a crossing, so the
%% probabilities rest on the background assumption (stated once, Sec. 4.3.1); the
%% rescaling to the measured near-trigger ratio is a sensitivity check.
%% eq:excl now labels the amplitude-fraction bound, the one the ledger uses.
\section{The control assumption, the hold-out and the recurrence bound}
\label{app:radius}

The chance expectation of \S\ref{sec:trigunif} predicts how often the star
should cross the trigger from how often control positions in its own field
do. It therefore rests on one assumption the searched data cannot test,
that the stellar position behaves like a control position. Below the trigger it does
not quite do so. In the band $\CeNearBandLo\le T<\CeNearBandHi$ the star
carries \CeNear{} times a control position's above-threshold cells
(\CeNearIntLo--\CeNearIntHi{} at \CeCovPct{} per cent, blocks resampled).
Rescaling the expectation by that ratio is a sensitivity check
(\S\ref{sec:trigunif}); the control positions' own rate, the more demanding
of the two, is kept for the count under test. For the \CeObs{}
unattributed crossings to stand above chance at \CeAlphaPct{} per cent, stars
would have to trigger at \CeRCrit{} times a control position's rate or less.

The stellar excess is not a primary-beam effect, since a noise rise toward the field
edge would put the stellar rank above one half, the opposite of what is
seen, and the probes' radial slope on the reserved blocks,
$\HoRadSlope\pm\HoRadSlopeSe$, is small and of the wrong sign. It appears
only in the 7\,m-array data, \BgArrSevenRatio{}
(\BgArrSevenRatioLo--\BgArrSevenRatioHi) against \BgArrTwelveRatio{}
(\BgArrTwelveRatioLo--\BgArrTwelveRatioHi) on the 12\,m array, and its cause
has not been identified.

A calibration error that varies with frequency, such as a bandpass residual
or a coarse system temperature interpolated across a telluric line
\citep{ALMATsys2018}, would act in the sense observed: it multiplies the
continuum, so it acts at the star and not on the empty sky around it. Such
a residual matters even if it is too small to make a crossing by itself. A small
error coupled to the continuum, added to noise already near the trigger,
would still change how often the star crosses it, and the searched data
cannot separate that from a real difference between star and controls. The
probabilities therefore rest on the stellar-background model
(\S\ref{sec:trigunif}).

\subsection{The reserved hold-out}
\label{sec:heldout}

Because the statistic and the mask were set on the searched sample,
\NEbHeldOut{} of \NEbProcessedAll{} blocks (\HoWin{} windows toward
\HoStars{} stars) were held out by an identifier rule,
$\mathrm{sha256}(\textsc{uid})_{[:8]} \bmod 5 = 0$. Two guards, which also
use identifiers alone, keep a block on a named list in the survey and ensure
that no system loses every block. Their independence of the statistic and the mask
is our assertion, which a reader cannot check. Since the rule hashes the
block and not the session, \HsSharePct{} per cent of them also share
calibrators, weather and setup with a searched block. They trigger in \StHoObs{} of
\StHoNFine{} windows against \StHoExpFine{} predicted, and the
\HsIndBlocks{} that share no session hold \HsIndObsWord{} trigger
\HsIndObsNoun{} against \HsIndExp, so they neither confirm nor contradict
the stellar ratio.

\subsection{The recurrence bound}
\label{sec:nullscatter}

A crossing is tested for recurrence by reading its predicted cell, one
channel at the carried drift, in the later blocks that cover it. With $S$
the discovery amplitude and $\bar{a}$ and $\bar{\sigma}=(\sum_i
\sigma_i^{-2})^{-1/2}$ the inverse-variance combination of the repeat
readings, the two signal-to-noise ratios are $T_{\rm pers}=S/\bar{\sigma}$
and $T_{\rm rep}=\bar{a}/\bar{\sigma}$. A discovery amplitude selected from
$10^{\TuCellsDecLo}$ to $10^{\TuCellsDecHi}$ cells at $T_\star\simeq5$ is
biased high, so it is not compared with the repeat directly. The bound is
instead the largest fraction of the discovery amplitude still allowed by the
repeat readings,
\begin{equation*}
f_{\rm rep}=\frac{T_{\rm rep}+5}{T_{\rm pers}},
\end{equation*}
and $T_{\rm rep}$ is centred on zero under the null, scattering by
\NuNullSd{} about \NuNullMean{} over \NuNullN{} single-cell repeat readings.
The reading is conditional. Given an assumed true amplitude, the bound
recovers a carrier with probability $\Phi(T_{\rm pers}-5)$ at the discovery
amplitude and $\Phi(5T_{\rm pers}/T_\star-5)$ at the nominal
threshold-equivalent amplitude $5S/T_\star$, and the \HjNPowerLo{} crossings
where these are smallest are reported as untested.
